\textbf{结论名称：}复合函数的单调性（同增异减）

\textbf{适用条件：}设内层函数 \(u=g(x)\) 在区间 \(I\) 上有单调性，外层函数 \(y=f(u)\) 在区间 \(U\) 上有单调性，且 \(g(I)\subseteq U\)，则复合函数 \(y=f(g(x))\) 在 \(I\) 上具有单调性。

\textbf{变量说明：}自变量 \(x\)，中间变量 \(u=g(x)\)，因变量 \(y=f(u)=f(g(x))\)；用 \(\nearrow\) 表示递增，\(\searrow\) 表示递减。

\textbf{核心公式：}
\[
\boxed{
\begin{array}{l}
f\nearrow,\;g\nearrow\;\Rightarrow\;f\circ g\nearrow;\qquad
f\searrow,\;g\searrow\;\Rightarrow\;f\circ g\nearrow\\[4pt]
f\nearrow,\;g\searrow\;\Rightarrow\;f\circ g\searrow;\qquad
f\searrow,\;g\nearrow\;\Rightarrow\;f\circ g\searrow
\end{array}
}
\]
简记为“同增异减”。

\textbf{使用场景：}快速判断复合型函数的单调区间，用于比较大小、解不等式、求参数范围等。

\textbf{简单例子：}已知 \(f(u)=\ln u\) 在 \((0,+\infty)\) 上递增（\(f\nearrow\)），\(g(x)=\frac{1}{x}\) 在 \((0,+\infty)\) 上递减（\(g\searrow\)），则 \(y=\ln\frac{1}{x}=-\ln x\) 在 \((0,+\infty)\) 上递减（异减）。验证：取 \(0<x_1<x_2\)，有 \(\frac{1}{x_1}>\frac{1}{x_2}\)，得 \(\ln\frac{1}{x_1}>\ln\frac{1}{x_2}\)，单调递减成立。